% ============================================================
% Academy Co-Scientist — FlowCept Provenance Data Schema
% ============================================================
\documentclass[11pt,a4paper]{article}

\usepackage[T1]{fontenc}
\usepackage[utf8]{inputenc}
\usepackage[margin=2.5cm]{geometry}
\usepackage{booktabs}
\usepackage{longtable}
\usepackage{array}
\usepackage{xcolor}
\usepackage{tabularx}
\usepackage{multirow}
\usepackage{hyperref}
\usepackage{parskip}
\usepackage{titlesec}
\usepackage{fancyhdr}
\usepackage{caption}
\usepackage{lmodern}
\usepackage{microtype}
\usepackage{forest}
\usepackage{listings}
\usepackage{enumitem}

\definecolor{fieldgray}{gray}{0.93}
\definecolor{subtypeblue}{RGB}{30,90,160}
\definecolor{optred}{RGB}{180,40,40}
\definecolor{condorange}{RGB}{200,100,0}

\hypersetup{
  colorlinks=true,
  linkcolor=subtypeblue,
  urlcolor=subtypeblue,
  pdftitle={Academy Co-Scientist Provenance Schema},
  pdfauthor={Academy Co-Scientist},
}

\pagestyle{fancy}
\fancyhf{}
\rhead{Academy Co-Scientist — Provenance Schema}
\lhead{}
\cfoot{\thepage}

\newcolumntype{L}[1]{>{\raggedright\arraybackslash}p{#1}}
\newcolumntype{C}[1]{>{\centering\arraybackslash}p{#1}}

% Shorthand for optional marker
\newcommand{\opt}{\textcolor{optred}{opt}}
\newcommand{\cond}{\textcolor{condorange}{cond}}
\newcommand{\req}{\textbf{req}}

% Subtype badge style
\newcommand{\subtype}[1]{\texttt{\textcolor{subtypeblue}{#1}}}

\captionsetup[table]{skip=4pt}

\title{\textbf{Academy Co-Scientist}\\[0.4em]
       \large FlowCept Provenance Data Schema\\[0.2em]
       \normalsize All record types emitted to the JSONL provenance buffer}
\author{}
\date{\today}

\begin{document}

\maketitle
\thispagestyle{fancy}
\tableofcontents
\newpage

% ============================================================
\section{Overview}
% ============================================================

All provenance records are serialised as JSON objects and appended line-by-line
to a \texttt{.jsonl} buffer file (one object per line).
Records are produced by four plugins:

\begin{itemize}[itemsep=2pt]
  \item \textbf{\texttt{FlowceptAcademyPlugin}} — patches \texttt{academy.runtime.Runtime}
        at the class level; captures every \texttt{@action}, \texttt{@loop},
        agent lifecycle event, and LLM/embedding call.
  \item \textbf{\texttt{FlowceptLangGraphPlugin}} — LangChain
        \texttt{BaseCallbackHandler} that captures graph runs, node executions,
        LLM calls, and tool calls.
  \item \textbf{\texttt{FlowceptCrewAIPlugin}} — combines CrewAI's native event
        bus (crew/task/agent lifecycle) with before/after LLM and tool hooks
        (richer data than the event bus alone).
  \item \textbf{\texttt{FlowceptAutoGenPlugin}} — wraps \texttt{team.run\_stream()}
        to capture AutoGen team runs and individual messages.
\end{itemize}

All plugins share the same \textbf{FlowCept} in-memory buffer and the same
top-level \texttt{workflow\_id} when combined via \texttt{from\_academy\_plugin()}.

\subsection*{Notation}
\begin{itemize}[itemsep=2pt]
  \item \req\ — field always present
  \item \opt\ — field present only when the relevant data is available
  \item \cond\ — field present conditionally (depends on \texttt{call\_type},
        status, or plugin configuration); explained in the table notes
\end{itemize}

% ============================================================
\section{Provenance Hierarchy}
% ============================================================

\begin{figure}[h!]
\centering
\begin{forest}
for tree={
  font=\small\ttfamily,
  grow'=0,
  draw,
  rounded corners,
  l sep=6mm,
  s sep=2mm,
  edge={-latex},
}
[campaign\_id (session)
  [workflow\_id (top-level run)
    [WorkflowObject: agent sub-wf \textrm{(Academy)}
      [academy\_lifecycle]
      [academy\_loop]
      [academy\_action
        [llm\_call]
      ]
    ]
    [WorkflowObject: graph/kickoff/run
      [langgraph\_graph / crewai\_crew / autogen\_run
        [langgraph\_node / crewai\_task / autogen\_message
          [llm\_call]
          [tool\_call]
        ]
        [crewai\_agent
          [llm\_call \textrm{(hook)}]
          [tool\_call \textrm{(hook)}]
        ]
      ]
    ]
  ]
]
\end{forest}
\caption{Provenance hierarchy. Arrows denote \texttt{parent\_task\_id} /
         \texttt{parent\_workflow\_id} linkage.}
\end{figure}

% ============================================================
\section{Common Envelope}
% ============================================================

Every record — regardless of plugin or subtype — carries the following fields.
They are injected by \texttt{intercept\_task()} and
\texttt{TaskObject.enrich\_task\_dict()}.

\begin{longtable}{L{3.5cm} L{2.5cm} C{1.2cm} L{7cm}}
\toprule
\textbf{Field} & \textbf{Type} & \textbf{P} & \textbf{Description} \\
\midrule
\endfirsthead
\multicolumn{4}{c}{\small\itshape (continued)} \\
\toprule
\textbf{Field} & \textbf{Type} & \textbf{P} & \textbf{Description} \\
\midrule
\endhead
\bottomrule
\endlastfoot

\texttt{task\_id}            & uuid string   & \req  & Unique identifier for this record \\
\rowcolor{fieldgray}
\texttt{workflow\_id}        & uuid string   & \req  & Top-level workflow for the run \\
\texttt{campaign\_id}        & uuid string   & \req  & FlowCept campaign (session-level grouping) \\
\rowcolor{fieldgray}
\texttt{subtype}             & string        & \req  & Record type discriminator — see Section~\ref{sec:subtypes} \\
\texttt{activity\_id}        & string        & \req  & Human-readable name: action name, node name, model ID, agent role, etc. \\
\rowcolor{fieldgray}
\texttt{group\_id}           & uuid string   & \opt  & Groups all tasks in one invocation (loop, graph run, kickoff, team run) \\
\texttt{parent\_task\_id}    & uuid string   & \opt  & Links child to parent record via \texttt{task\_id} \\
\rowcolor{fieldgray}
\texttt{started\_at}         & float         & \req  & Wall-clock start time (Unix epoch seconds) \\
\texttt{ended\_at}           & float         & \req  & Wall-clock end time (Unix epoch seconds) \\
\rowcolor{fieldgray}
\texttt{status}              & string enum   & \req  & \texttt{"FINISHED"} or \texttt{"ERROR"} \\
\texttt{stderr}              & string        & \opt  & Error message; present only when \texttt{status = ERROR} \\
\rowcolor{fieldgray}
\texttt{used}                & object        & \opt  & Input / request data — structure varies by subtype \\
\texttt{generated}           & object        & \opt  & Output / response data — structure varies by subtype \\
\rowcolor{fieldgray}
\texttt{custom\_metadata}    & object        & \opt  & Framework-specific context \\
\texttt{telemetry\_at\_start}& object        & \opt  & CPU/memory snapshot at start (when telemetry is enabled) \\
\rowcolor{fieldgray}
\texttt{telemetry\_at\_end}  & object        & \opt  & CPU/memory snapshot at end \\
\texttt{node\_name}          & string        & \opt  & Compute node name — from \texttt{enrich\_task\_dict()} \\
\rowcolor{fieldgray}
\texttt{login\_name}         & string        & \opt  & OS user name \\
\texttt{hostname}            & string        & \opt  & Full hostname \\
\rowcolor{fieldgray}
\texttt{public\_ip}          & string        & \opt  & Public IP of the node \\
\texttt{private\_ip}         & string        & \opt  & Private IP of the node \\

\caption{Common record envelope — all subtypes.}
\label{tab:common}
\end{longtable}

% ============================================================
\section{Record Types by Subtype}
\label{sec:subtypes}
% ============================================================

% ------------------------------------------------------------
\subsection{Academy Plugin (\texttt{flowcept\_plugin.py})}
% ------------------------------------------------------------

\subsubsection{\subtype{academy\_action}}

Emitted once per \texttt{@action} dispatch by the patched
\texttt{Runtime.action()}. In addition to the common envelope:

\begin{longtable}{L{4cm} L{2.5cm} C{1.2cm} L{6.5cm}}
\toprule
\textbf{Field} & \textbf{Type} & \textbf{P} & \textbf{Description} \\
\midrule
\endfirsthead
\toprule
\textbf{Field} & \textbf{Type} & \textbf{P} & \textbf{Description} \\
\midrule
\endhead
\bottomrule
\endlastfoot

\texttt{agent\_id}
  & string & \req
  & Real Academy agent identifier (\texttt{AgentId<hex>}) \\
\rowcolor{fieldgray}
\multicolumn{4}{l}{\textit{\textbf{used} object}} \\
\texttt{used.args}
  & any & \req & Positional arguments passed to the action \\
\rowcolor{fieldgray}
\texttt{used.kwargs}
  & object & \req & Keyword arguments passed to the action \\
\multicolumn{4}{l}{\textit{\textbf{generated} object}} \\
\rowcolor{fieldgray}
\texttt{generated}
  & any & \opt & Return value of the action (absent on error) \\
\multicolumn{4}{l}{\textit{\textbf{custom\_metadata} object}} \\
\rowcolor{fieldgray}
\texttt{agent\_type}
  & string & \req & Agent class name (e.g.\ \texttt{TeamCoordinatorAgent}) \\
\texttt{source\_id}
  & string & \req & Caller identity: \texttt{AgentId<hex>} or \texttt{UserId<hex>} \\
\rowcolor{fieldgray}
\texttt{cross\_agent\_call}
  & bool & \req & \texttt{true} when the caller is another agent \\
\texttt{source\_agent\_id}
  & string & \opt & Set when \texttt{cross\_agent\_call} is true \\
\rowcolor{fieldgray}
\texttt{source\_workflow\_id}
  & uuid string & \opt & Sub-workflow ID of the calling agent \\

\caption{\subtype{academy\_action} record fields.}
\end{longtable}

\subsubsection{\subtype{academy\_loop}}

Emitted at loop \textbf{start}, \textbf{exit}, and \textbf{error} by the patched
\texttt{Runtime.\_execute\_loop()}. All three events share the same
\texttt{task\_id} and \texttt{group\_id}.

\begin{longtable}{L{4cm} L{2.5cm} C{1.2cm} L{6.5cm}}
\toprule
\textbf{Field} & \textbf{Type} & \textbf{P} & \textbf{Description} \\
\midrule
\endfirsthead
\toprule \textbf{Field} & \textbf{Type} & \textbf{P} & \textbf{Description} \\
\midrule
\endhead
\bottomrule
\endlastfoot

\texttt{agent\_id} & string & \req & \texttt{AgentId<hex>} \\
\rowcolor{fieldgray}
\multicolumn{4}{l}{\textit{\textbf{custom\_metadata} object}} \\
\texttt{agent\_type} & string & \req & Agent class name \\
\rowcolor{fieldgray}
\texttt{loop\_event}
  & string enum & \req & \texttt{"start"} | \texttt{"exit"} | \texttt{"error"} \\

\caption{\subtype{academy\_loop} record fields.}
\end{longtable}

\subsubsection{\subtype{academy\_lifecycle}}

One record each for \texttt{agent\_startup} and \texttt{agent\_shutdown}.
\texttt{activity\_id} is the event name.

\begin{longtable}{L{4cm} L{2.5cm} C{1.2cm} L{6.5cm}}
\toprule
\textbf{Field} & \textbf{Type} & \textbf{P} & \textbf{Description} \\
\midrule
\endfirsthead
\toprule \textbf{Field} & \textbf{Type} & \textbf{P} & \textbf{Description} \\
\midrule
\endhead
\bottomrule
\endlastfoot

\texttt{agent\_id} & string & \req & \texttt{AgentId<hex>} \\
\rowcolor{fieldgray}
\multicolumn{4}{l}{\textit{\textbf{custom\_metadata} object}} \\
\texttt{agent\_type} & string & \req & Agent class name \\

\caption{\subtype{academy\_lifecycle} record fields.}
\end{longtable}

% ========================
\subsubsection{\subtype{llm\_call} (Academy)}
\label{sec:llm-academy}
% ========================

Emitted inside \texttt{\_process\_llm\_call()} via the LLM hook registered by
\texttt{FlowceptAcademyPlugin}. The \texttt{used} and \texttt{generated}
sub-objects vary by \texttt{call\_type}.

\paragraph{Common \& identity fields}

\begin{longtable}{L{4cm} L{2.5cm} C{1.2cm} L{6.5cm}}
\toprule
\textbf{Field} & \textbf{Type} & \textbf{P} & \textbf{Description} \\
\midrule
\endfirsthead
\toprule \textbf{Field} & \textbf{Type} & \textbf{P} & \textbf{Description} \\
\midrule
\endhead
\bottomrule
\endlastfoot

\texttt{agent\_id}
  & string & \req
  & \texttt{AgentId<hex>} from \texttt{\_current\_academy\_agent\_id} ContextVar;
    falls back to agent class name \\
\rowcolor{fieldgray}
\texttt{parent\_task\_id}
  & uuid string & \opt
  & ID of the enclosing \subtype{academy\_action} task from
    \texttt{\_current\_action\_task\_id} ContextVar \\

\caption{\subtype{llm\_call} (Academy) — identity fields.}
\end{longtable}

\paragraph{\texttt{used} object — request side}

\begin{longtable}{L{4cm} L{2.5cm} C{1.2cm} L{6.5cm}}
\toprule
\textbf{Field} & \textbf{Type} & \textbf{P} & \textbf{Description} \\
\midrule
\endfirsthead
\toprule \textbf{Field} & \textbf{Type} & \textbf{P} & \textbf{Description} \\
\midrule
\endhead
\bottomrule
\endlastfoot

\texttt{model}
  & string & \req & Requested model identifier \\
\rowcolor{fieldgray}
\texttt{model\_used}
  & string & \req & Model actually used (may differ due to routing) \\
\texttt{call\_type}
  & string & \req
  & Logical call type from \texttt{ctx["call\_type"]}; one of
    \texttt{parsed\_json\_result}, \texttt{chat\_completion},
    \texttt{embed\_result\_local}, \texttt{embed\_result\_openai} \\
\rowcolor{fieldgray}
\texttt{temperature\_requested}
  & float & \opt & Temperature value passed to the call \\
\texttt{temperature\_sent}
  & float & \opt & Temperature value actually sent to the API \\
\rowcolor{fieldgray}
\texttt{temperature\_suppressed}
  & bool & \req
  & \texttt{true} when the model ignores temperature
    (e.g.\ o1, o4-mini reasoning models) \\
\texttt{temperature\_null\_reason}
  & string & \cond
  & Present when temperature is null; one of
    \texttt{"reasoning\_model\_does\_not\_accept\_temperature"},
    \texttt{"not\_applicable\_for\_embeddings"},
    \texttt{"not\_set"} \\
\rowcolor{fieldgray}
\texttt{messages}
  & list & \opt & Full message list when present in the payload \\
\texttt{system\_instructions}
  & string & \opt & System prompt \\
\rowcolor{fieldgray}
\texttt{user\_prompt}
  & string & \opt & User message \\
\texttt{schema\_hint}
  & any & \opt & JSON schema for structured output \\
\rowcolor{fieldgray}
\texttt{agent\_class}
  & string & \req & Agent class name (e.g.\ \texttt{HypothesisRefinerAgent}) \\
\texttt{academy\_agent\_id}
  & string & \opt & Raw \texttt{AgentId<hex>} string \\
\rowcolor{fieldgray}
\texttt{ctx\_<key>}
  & any & \opt
  & Every field from the caller's \texttt{context=\{\}} dict is hoisted
    with a \texttt{ctx\_} prefix
    (e.g.\ \texttt{ctx\_hyp\_id}, \texttt{ctx\_team},
    \texttt{ctx\_instance\_id}, \texttt{ctx\_call\_type}) \\

\caption{\subtype{llm\_call} (Academy) — \texttt{used} object.}
\end{longtable}

\paragraph{\texttt{generated} object — response side, by \texttt{call\_type}}

\noindent\textbf{(a)~\texttt{parsed\_json\_result}} — structured LLM call via
\texttt{\_call\_llm\_json}:

\begin{longtable}{L{4cm} L{2.5cm} C{1.2cm} L{6.5cm}}
\toprule
\textbf{Field} & \textbf{Type} & \textbf{P} & \textbf{Description} \\
\midrule
\endfirsthead
\toprule \textbf{Field} & \textbf{Type} & \textbf{P} & \textbf{Description} \\
\midrule
\endhead
\bottomrule
\endlastfoot

\texttt{finish\_reason}      & string & \opt & API finish reason \\
\rowcolor{fieldgray}
\texttt{prompt\_tokens}      & int    & \opt & Token usage \\
\texttt{completion\_tokens}  & int    & \opt & Token usage \\
\rowcolor{fieldgray}
\texttt{total\_tokens}       & int    & \opt & Total token usage \\
\texttt{parsed\_response}    & object & \req & Full structured dict returned by the LLM \\
\rowcolor{fieldgray}
\texttt{fallback\_used}      & bool   & \req & Whether regex fallback was needed to extract JSON \\
\texttt{raw\_response\_text} & string & \req & Raw text before JSON parsing \\
\rowcolor{fieldgray}
\texttt{<key>}
  & any & \cond
  & \textbf{Every top-level key of \texttt{parsed\_response} is also hoisted
    directly} into \texttt{generated} for direct queryability
    (e.g.\ \texttt{score}, \texttt{title}, \texttt{description},
    \texttt{recommendation}, \texttt{weaknesses}) \\

\caption{\subtype{llm\_call} (Academy) — \texttt{generated} for \texttt{parsed\_json\_result}.}
\end{longtable}

\noindent\textbf{(b)~\texttt{chat\_completion}} — raw text call:

\begin{longtable}{L{4cm} L{2.5cm} C{1.2cm} L{6.5cm}}
\toprule
\textbf{Field} & \textbf{Type} & \textbf{P} & \textbf{Description} \\
\midrule
\endfirsthead
\toprule \textbf{Field} & \textbf{Type} & \textbf{P} & \textbf{Description} \\
\midrule
\endhead
\bottomrule
\endlastfoot

\texttt{finish\_reason}     & string & \opt & API finish reason \\
\rowcolor{fieldgray}
\texttt{prompt\_tokens}     & int    & \opt & Token usage \\
\texttt{completion\_tokens} & int    & \opt & \\
\rowcolor{fieldgray}
\texttt{total\_tokens}      & int    & \opt & \\
\texttt{response\_text}     & string & \req & Full raw text response \\
\rowcolor{fieldgray}
\texttt{parsed\_response}   & object & \opt & Inline JSON parsed from the text, if valid \\
\texttt{<key>}              & any    & \cond
  & Top-level keys from inline JSON hoisted directly \\

\caption{\subtype{llm\_call} (Academy) — \texttt{generated} for \texttt{chat\_completion}.}
\end{longtable}

\noindent\textbf{(c)~\texttt{embed\_result\_local} / \texttt{embed\_result\_openai}} — embedding calls:

\begin{longtable}{L{4cm} L{2.5cm} C{1.2cm} L{6.5cm}}
\toprule
\textbf{Field} & \textbf{Type} & \textbf{P} & \textbf{Description} \\
\midrule
\endfirsthead
\toprule \textbf{Field} & \textbf{Type} & \textbf{P} & \textbf{Description} \\
\midrule
\endhead
\bottomrule
\endlastfoot

\texttt{num\_vectors} & int & \opt & Number of embedding vectors returned \\
\rowcolor{fieldgray}
\texttt{embed\_dim}   & int & \opt & Dimensionality of each embedding vector \\

\caption{\subtype{llm\_call} (Academy) — \texttt{generated} for embedding calls.}
\end{longtable}

\noindent\textbf{Error case} (any \texttt{call\_type} when \texttt{status = ERROR}):

\vspace{0.3em}
\texttt{generated.error} \textit{(string)} — \texttt{str(exception)}.

\paragraph{\texttt{custom\_metadata} object}

\begin{longtable}{L{4cm} L{2.5cm} C{1.2cm} L{6.5cm}}
\toprule
\textbf{Field} & \textbf{Type} & \textbf{P} & \textbf{Description} \\
\midrule
\endfirsthead
\toprule \textbf{Field} & \textbf{Type} & \textbf{P} & \textbf{Description} \\
\midrule
\endhead
\bottomrule
\endlastfoot

\texttt{llm\_call\_type} & string & \req & Raw \texttt{call\_type} value \\
\rowcolor{fieldgray}
\texttt{hyp\_id}         & string & \opt & Hypothesis ID from \texttt{ctx["hyp\_id"]} \\

\caption{\subtype{llm\_call} (Academy) — \texttt{custom\_metadata}.}
\end{longtable}

% ------------------------------------------------------------
\subsection{LangGraph Plugin (\texttt{flowcept\_langgraph\_plugin.py})}
% ------------------------------------------------------------

\subsubsection{\subtype{langgraph\_graph}}

One record per \texttt{graph.invoke()} / \texttt{graph.ainvoke()} call.

\begin{longtable}{L{4cm} L{2.5cm} C{1.2cm} L{6.5cm}}
\toprule
\textbf{Field} & \textbf{Type} & \textbf{P} & \textbf{Description} \\
\midrule
\endfirsthead
\toprule \textbf{Field} & \textbf{Type} & \textbf{P} & \textbf{Description} \\
\midrule
\endhead
\bottomrule
\endlastfoot

\texttt{used.inputs}           & object & \req & Graph input state \\
\rowcolor{fieldgray}
\texttt{generated.outputs}     & object & \req & Graph output state \\
\texttt{custom\_metadata.tags} & list   & \req & LangChain tags \\
\rowcolor{fieldgray}
\texttt{custom\_metadata.graph\_name}      & string & \req & Graph name \\
\texttt{custom\_metadata.source\_agent\_id}& string & \opt
  & Academy \texttt{AgentId} that produced the input
    (from \texttt{state["\_source\_agent\_id"]}) \\

\caption{\subtype{langgraph\_graph} record fields.}
\end{longtable}

\subsubsection{\subtype{langgraph\_node}}

One record per node execution within a graph run.

\begin{longtable}{L{4cm} L{2.5cm} C{1.2cm} L{6.5cm}}
\toprule
\textbf{Field} & \textbf{Type} & \textbf{P} & \textbf{Description} \\
\midrule
\endfirsthead
\toprule \textbf{Field} & \textbf{Type} & \textbf{P} & \textbf{Description} \\
\midrule
\endhead
\bottomrule
\endlastfoot

\texttt{used.inputs}            & object & \req & Node input state \\
\rowcolor{fieldgray}
\texttt{generated.outputs}      & object & \req & Node output state \\
\texttt{custom\_metadata.tags}  & list   & \req & LangChain tags \\
\rowcolor{fieldgray}
\texttt{custom\_metadata.node\_name}       & string & \req & Node name \\
\texttt{custom\_metadata.source\_agent\_id}& string & \opt & Cross-framework linkage \\
\rowcolor{fieldgray}
\texttt{parent\_task\_id}       & uuid   & \opt & Enclosing \subtype{langgraph\_graph} \texttt{task\_id} \\

\caption{\subtype{langgraph\_node} record fields.}
\end{longtable}

\subsubsection{\subtype{llm\_call} (LangGraph)}

Emitted by \texttt{on\_llm\_start} + \texttt{on\_llm\_end} (completion models)
or \texttt{on\_chat\_model\_start} + \texttt{on\_llm\_end} (chat models).

\begin{longtable}{L{4cm} L{2.5cm} C{1.2cm} L{6.5cm}}
\toprule
\textbf{Field} & \textbf{Type} & \textbf{P} & \textbf{Description} \\
\midrule
\endfirsthead
\toprule \textbf{Field} & \textbf{Type} & \textbf{P} & \textbf{Description} \\
\midrule
\endhead
\bottomrule
\endlastfoot

\multicolumn{4}{l}{\textit{\textbf{used} object}} \\
\rowcolor{fieldgray}
\texttt{prompts}  & list[string] & \cond & Raw prompt strings (\texttt{on\_llm\_start} path) \\
\texttt{messages} & list[list[string]] & \cond
  & Serialised message content arrays (\texttt{on\_chat\_model\_start} path) \\
\rowcolor{fieldgray}
\texttt{model}    & string & \req & Model name from \texttt{serialized.kwargs} \\
\multicolumn{4}{l}{\textit{\textbf{generated} object}} \\
\rowcolor{fieldgray}
\texttt{text}               & string & \req & First generation text from \texttt{LLMResult} \\
\texttt{prompt\_tokens}     & int    & \opt & From \texttt{llm\_output.token\_usage} \\
\rowcolor{fieldgray}
\texttt{completion\_tokens} & int    & \opt & \\
\texttt{total\_tokens}      & int    & \opt & \\
\rowcolor{fieldgray}
\texttt{model}              & string & \req & Resolved model from \texttt{llm\_output.model\_name} \\
\multicolumn{4}{l}{\textit{\textbf{custom\_metadata} object}} \\
\rowcolor{fieldgray}
\texttt{model}              & string & \req & Model name \\
\multicolumn{4}{l}{\textit{Linkage}} \\
\rowcolor{fieldgray}
\texttt{parent\_task\_id}   & uuid   & \opt & Enclosing \subtype{langgraph\_node} \texttt{task\_id} \\

\caption{\subtype{llm\_call} (LangGraph) record fields.}
\end{longtable}

\subsubsection{\subtype{tool\_call} (LangGraph)}

Emitted by \texttt{on\_tool\_start} + \texttt{on\_tool\_end}.

\begin{longtable}{L{4cm} L{2.5cm} C{1.2cm} L{6.5cm}}
\toprule
\textbf{Field} & \textbf{Type} & \textbf{P} & \textbf{Description} \\
\midrule
\endfirsthead
\toprule \textbf{Field} & \textbf{Type} & \textbf{P} & \textbf{Description} \\
\midrule
\endhead
\bottomrule
\endlastfoot

\texttt{used.input}              & string & \req & Tool input string \\
\rowcolor{fieldgray}
\texttt{generated.output}        & string & \req & Tool output string \\
\texttt{custom\_metadata.tool\_name} & string & \req & From \texttt{serialized["name"]} \\
\rowcolor{fieldgray}
\texttt{parent\_task\_id}        & uuid   & \opt & Enclosing \subtype{langgraph\_node} \texttt{task\_id} \\

\caption{\subtype{tool\_call} (LangGraph) record fields.}
\end{longtable}

% ------------------------------------------------------------
\subsection{CrewAI Plugin (\texttt{flowcept\_crewai\_plugin.py})}
% ------------------------------------------------------------

\subsubsection{\subtype{crewai\_crew}}

One record per \texttt{crew.kickoff()} call; emitted via the event bus.

\begin{longtable}{L{4cm} L{2.5cm} C{1.2cm} L{6.5cm}}
\toprule
\textbf{Field} & \textbf{Type} & \textbf{P} & \textbf{Description} \\
\midrule
\endfirsthead
\toprule \textbf{Field} & \textbf{Type} & \textbf{P} & \textbf{Description} \\
\midrule
\endhead
\bottomrule
\endlastfoot

\texttt{used.inputs}                & object & \req & \texttt{kickoff(inputs=...)} dict \\
\rowcolor{fieldgray}
\texttt{used.crew\_name}            & string & \req & Crew name \\
\texttt{generated.output}           & any    & \req & Crew output \\
\rowcolor{fieldgray}
\texttt{generated.total\_tokens}    & int    & \opt & Aggregate token count if available \\
\texttt{custom\_metadata.crew\_name}& string & \req & Crew name \\
\rowcolor{fieldgray}
\texttt{custom\_metadata.framework} & string & \req & \texttt{"crewai"} \\

\caption{\subtype{crewai\_crew} record fields.}
\end{longtable}

\subsubsection{\subtype{crewai\_task}}

One record per CrewAI task execution; emitted via the event bus.

\begin{longtable}{L{4cm} L{2.5cm} C{1.2cm} L{6.5cm}}
\toprule
\textbf{Field} & \textbf{Type} & \textbf{P} & \textbf{Description} \\
\midrule
\endfirsthead
\toprule \textbf{Field} & \textbf{Type} & \textbf{P} & \textbf{Description} \\
\midrule
\endhead
\bottomrule
\endlastfoot

\texttt{used.context}               & object & \opt & Task context dict \\
\rowcolor{fieldgray}
\texttt{used.task}                  & object & \opt & Serialised task object \\
\texttt{generated.output}           & any    & \req & Task output \\
\rowcolor{fieldgray}
\texttt{custom\_metadata.task\_name}& string & \req & Task name \\
\texttt{custom\_metadata.agent\_role}& string & \req & Role of the assigned agent \\
\rowcolor{fieldgray}
\texttt{custom\_metadata.framework} & string & \req & \texttt{"crewai"} \\

\caption{\subtype{crewai\_task} record fields.}
\end{longtable}

\subsubsection{\subtype{crewai\_agent}}

One record per agent execution within a task; emitted via the event bus.

\begin{longtable}{L{4cm} L{2.5cm} C{1.2cm} L{6.5cm}}
\toprule
\textbf{Field} & \textbf{Type} & \textbf{P} & \textbf{Description} \\
\midrule
\endfirsthead
\toprule \textbf{Field} & \textbf{Type} & \textbf{P} & \textbf{Description} \\
\midrule
\endhead
\bottomrule
\endlastfoot

\texttt{used.task\_prompt}           & string & \req & Full task prompt seen by the agent \\
\rowcolor{fieldgray}
\texttt{used.tools}                  & list   & \req & Tool names available to the agent \\
\texttt{generated.output}            & any    & \req & Agent execution output \\
\rowcolor{fieldgray}
\texttt{custom\_metadata.agent\_role}& string & \req & Agent role \\
\texttt{custom\_metadata.framework}  & string & \req & \texttt{"crewai"} \\
\rowcolor{fieldgray}
\texttt{parent\_task\_id}            & uuid   & \opt & Enclosing \subtype{crewai\_task} \texttt{task\_id} \\

\caption{\subtype{crewai\_agent} record fields.}
\end{longtable}

\subsubsection{\subtype{llm\_call} (CrewAI hook)}
\label{sec:llm-crewai}

Emitted by \texttt{before\_llm\_call} + \texttt{after\_llm\_call} hooks.
These records are \textbf{richer} than event-bus LLM events because the hooks
receive the full agent and task context objects.

\begin{longtable}{L{4cm} L{2.5cm} C{1.2cm} L{6.5cm}}
\toprule
\textbf{Field} & \textbf{Type} & \textbf{P} & \textbf{Description} \\
\midrule
\endfirsthead
\toprule \textbf{Field} & \textbf{Type} & \textbf{P} & \textbf{Description} \\
\midrule
\endhead
\bottomrule
\endlastfoot

\multicolumn{4}{l}{\textit{\textbf{used} object}} \\
\rowcolor{fieldgray}
\texttt{messages}
  & list[\{role, content\}] & \req
  & Full message list passed to the LLM with \texttt{role} and \texttt{content}
    for each turn \\
\texttt{model}
  & string & \req & From \texttt{context.llm.model} or \texttt{context.llm.model\_name} \\
\rowcolor{fieldgray}
\texttt{iterations}
  & int & \req & Current ReAct loop iteration count \\
\texttt{agent\_role}
  & string & \req & \texttt{context.agent.role} \\
\rowcolor{fieldgray}
\texttt{agent\_goal}
  & string & \opt & \texttt{context.agent.goal} \\
\texttt{task\_description}
  & string & \opt & \texttt{context.task.description} \\
\multicolumn{4}{l}{\textit{\textbf{generated} object}} \\
\rowcolor{fieldgray}
\texttt{response}
  & string (max 4000 chars) & \req
  & Full LLM response text from \texttt{context.response} \\
\texttt{model}
  & string & \req & Model that was used \\
\multicolumn{4}{l}{\textit{\textbf{custom\_metadata} object}} \\
\rowcolor{fieldgray}
\texttt{model}       & string & \req & Model name \\
\texttt{agent\_role} & string & \req & Agent role \\
\rowcolor{fieldgray}
\texttt{framework}   & string & \req & \texttt{"crewai"} \\
\texttt{source}      & string & \req & \texttt{"llm\_hook"} — distinguishes from event-bus records \\
\rowcolor{fieldgray}
\multicolumn{4}{l}{\textit{Linkage}} \\
\texttt{parent\_task\_id}
  & uuid & \opt & Best-effort lookup of enclosing \subtype{crewai\_agent} \texttt{task\_id} \\

\caption{\subtype{llm\_call} (CrewAI hook) record fields.}
\end{longtable}

\subsubsection{\subtype{tool\_call} (CrewAI hook)}
\label{sec:tool-crewai}

Emitted by \texttt{before\_tool\_call} + \texttt{after\_tool\_call} hooks.
Tool inputs are \textbf{typed dicts} (not flat strings as in LangGraph).

\begin{longtable}{L{4cm} L{2.5cm} C{1.2cm} L{6.5cm}}
\toprule
\textbf{Field} & \textbf{Type} & \textbf{P} & \textbf{Description} \\
\midrule
\endfirsthead
\toprule \textbf{Field} & \textbf{Type} & \textbf{P} & \textbf{Description} \\
\midrule
\endhead
\bottomrule
\endlastfoot

\multicolumn{4}{l}{\textit{\textbf{used} object}} \\
\rowcolor{fieldgray}
\texttt{input}
  & object & \req
  & \textbf{Typed} tool input dict from \texttt{context.tool\_input}
    (richer than the LangGraph string) \\
\texttt{agent\_role}
  & string & \req & Agent role \\
\rowcolor{fieldgray}
\texttt{task\_description}
  & string & \opt & \texttt{context.task.description} \\
\multicolumn{4}{l}{\textit{\textbf{generated} object}} \\
\rowcolor{fieldgray}
\texttt{output}
  & string (max 2000 chars) & \req & Tool result from \texttt{context.tool\_result} \\
\multicolumn{4}{l}{\textit{\textbf{custom\_metadata} object}} \\
\rowcolor{fieldgray}
\texttt{tool\_name}  & string & \req & \texttt{context.tool\_name} \\
\texttt{agent\_role} & string & \req & Agent role \\
\rowcolor{fieldgray}
\texttt{framework}   & string & \req & \texttt{"crewai"} \\
\texttt{source}      & string & \req & \texttt{"tool\_hook"} \\
\rowcolor{fieldgray}
\multicolumn{4}{l}{\textit{Linkage}} \\
\texttt{parent\_task\_id}
  & uuid & \opt & Best-effort lookup of enclosing \subtype{crewai\_agent} \texttt{task\_id} \\

\caption{\subtype{tool\_call} (CrewAI hook) record fields.}
\end{longtable}

% ------------------------------------------------------------
\subsection{AutoGen Plugin (\texttt{flowcept\_autogen\_plugin.py})}
% ------------------------------------------------------------

\subsubsection{\subtype{autogen\_run}}

One record per \texttt{team.run()} call, emitted after the stream finishes.

\begin{longtable}{L{4cm} L{2.5cm} C{1.2cm} L{6.5cm}}
\toprule
\textbf{Field} & \textbf{Type} & \textbf{P} & \textbf{Description} \\
\midrule
\endfirsthead
\toprule \textbf{Field} & \textbf{Type} & \textbf{P} & \textbf{Description} \\
\midrule
\endhead
\bottomrule
\endlastfoot

\texttt{used.task}
  & string & \req & Task / initial message sent to the team \\
\rowcolor{fieldgray}
\texttt{generated.stop\_reason}
  & string & \req & Stop reason from \texttt{TaskResult} \\
\texttt{generated.message\_count}
  & int    & \req & Total number of messages captured \\
\rowcolor{fieldgray}
\texttt{generated.last\_message}
  & string & \req & Content of the last message (first 200 chars) \\
\texttt{generated.messages}
  & list[\{source, content\}] & \req
  & Summary list of all captured messages \\
\rowcolor{fieldgray}
\texttt{custom\_metadata.team\_name}       & string & \req & Team name \\
\texttt{custom\_metadata.framework}        & string & \req & \texttt{"autogen"} \\
\rowcolor{fieldgray}
\texttt{custom\_metadata.source\_agent\_id}& string & \opt
  & Academy \texttt{AgentId} that produced the task input \\

\caption{\subtype{autogen\_run} record fields.}
\end{longtable}

\subsubsection{\subtype{autogen\_message}}

One record per message yielded by \texttt{run\_stream()}.

\begin{longtable}{L{4cm} L{2.5cm} C{1.2cm} L{6.5cm}}
\toprule
\textbf{Field} & \textbf{Type} & \textbf{P} & \textbf{Description} \\
\midrule
\endfirsthead
\toprule \textbf{Field} & \textbf{Type} & \textbf{P} & \textbf{Description} \\
\midrule
\endhead
\bottomrule
\endlastfoot

\texttt{used.task}    & string & \req & Original task prompt \\
\rowcolor{fieldgray}
\texttt{used.agent}   & string & \req & Source agent name \\
\texttt{generated.content}
  & string | object & \req & Message content \\
\rowcolor{fieldgray}
\texttt{generated.message\_type}
  & string & \req & Python class name of the message object \\
\texttt{custom\_metadata.agent\_name}   & string & \req & Source agent name \\
\rowcolor{fieldgray}
\texttt{custom\_metadata.message\_type} & string & \req & Message class name \\
\texttt{custom\_metadata.framework}     & string & \req & \texttt{"autogen"} \\
\rowcolor{fieldgray}
\texttt{parent\_task\_id}
  & uuid & \req & Enclosing \subtype{autogen\_run} \texttt{task\_id} \\

\caption{\subtype{autogen\_message} record fields.}
\end{longtable}

% ============================================================
\section{Cross-Framework Comparison}
% ============================================================

Table~\ref{tab:comparison} summarises the capabilities of each \subtype{llm\_call}
variant and the \subtype{tool\_call} variants.

\begin{table}[h!]
\centering
\small
\begin{tabularx}{\textwidth}{L{4.5cm} C{2cm} C{2.5cm} C{2.5cm}}
\toprule
\textbf{Capability}
  & \textbf{Academy}
  & \textbf{LangGraph}
  & \textbf{CrewAI (hook)} \\
\midrule
Full message list captured  & yes & yes (chat path) & yes (with roles) \\
\rowcolor{fieldgray}
Token counts                & yes & yes             & no \\
Temperature details         & yes & no              & no \\
\rowcolor{fieldgray}
Agent context               & \texttt{ctx\_*} prefix & no & \texttt{agent\_role}, \texttt{agent\_goal} \\
Task / hypothesis context   & \texttt{ctx\_hyp\_id}  & no & \texttt{task\_description} \\
\rowcolor{fieldgray}
ReAct iteration count       & no  & no              & yes (\texttt{iterations}) \\
Structured output parsed    & yes (hoisted) & partial (inline JSON) & no \\
\rowcolor{fieldgray}
Linked to parent task       & \subtype{academy\_action} & \subtype{langgraph\_node} & \subtype{crewai\_agent} \\
\midrule
\textbf{Tool input type}    & n/a & string          & typed dict \\
\rowcolor{fieldgray}
Tool output captured        & n/a & yes             & yes \\
Agent/task context in tool  & n/a & no              & yes \\
\bottomrule
\end{tabularx}
\caption{Comparison of \subtype{llm\_call} and \subtype{tool\_call} across plugins.}
\label{tab:comparison}
\end{table}

% ============================================================
\section{Performance Overhead CSV}
% ============================================================

Each plugin optionally writes a per-call overhead CSV to
\texttt{provenance\_perf\_<workflow\_id>.csv}.

\begin{longtable}{L{3.5cm} L{2.5cm} L{7.5cm}}
\toprule
\textbf{Column} & \textbf{Type} & \textbf{Description} \\
\midrule
\endfirsthead
\toprule \textbf{Column} & \textbf{Type} & \textbf{Description} \\
\midrule
\endhead
\bottomrule
\endlastfoot

\texttt{timestamp\_utc} & ISO-8601 string & UTC time of the event \\
\rowcolor{fieldgray}
\texttt{workflow\_id}   & uuid string & Workflow this event belongs to \\
\texttt{category}       & string
  & One of: \texttt{action\_emit}, \texttt{loop\_emit},
    \texttt{lifecycle\_emit}, \texttt{llm\_hook},
    \texttt{intercept\_task}, \texttt{flush} \\
\rowcolor{fieldgray}
\texttt{elapsed\_us}    & float & Wall-clock overhead in microseconds \\

\caption{Provenance overhead CSV columns.}
\end{longtable}

\end{document}
